\chapter{Pop: Turning Feeling into a Hook}
\label{ch:genre-pop-deep}

\classquote{情发于声，声成文谓之音。}{《毛诗序》}{Emotion issues in sound; when sound takes form, it becomes music.}

\section{What pop is: four things behind ``catchy''}
\label{sec:pop-what}

\term{pop}{pop} is not a synonym for ``simple music'', and it is not what happens when you turn the volume up and fill every track. It is closer to a contract with the listener: in the first thirty seconds, give them a door they can remember; in three minutes, give them one complete emotional ride. The melody can be complex or five notes long; the instrumentation can be a full band or just a voice, a drum and a bass. What makes it pop is whether it says clearly what it wants you to hear.

Unpack ``catchy'' and at least four handles appear:

\begin{longtable}{p{0.16\textwidth}p{0.32\textwidth}p{0.38\textwidth}}
\toprule
Handle & What it feels like & The writing question \\
\midrule
\endfirsthead
\toprule
Handle & What it feels like & The writing question \\
\midrule
\endhead
Hook & A phrase, riff, timbre or groove that stays in the ear & What is the smallest memorable unit? If you remove it, does the song still stand? \\
Groove & The body understands before the mind does; a foot or hand starts moving & Which beat makes people nod? Is that accent clear in every bar? \\
Contrast & Verse and chorus differ in weight, light, density or colour & Can a listener tell ``this is the chorus'' within ten seconds? \\
Payoff & The expectation planted earlier is actually caught & Does the chorus catch the breath the pre-chorus lifted? \\
\bottomrule
\end{longtable}

The opposite of pop is not ``tasteful''; it is ``no choice''. A song where every section is equally full, every layer equally loud and every line equally long gives the listener nothing to allocate attention to. Professional pop often sounds restrained: it knows which layer should step back, which beat should stay empty, and which line should be saved for the next section.

\section{From gramophone to playlist: how pop got here}
\label{sec:pop-history}

Pop did not appear fully formed from one style. It absorbed dance music, folk song, jazz, rock, soul, electronics and hip-hop, then reorganised what it had borrowed into the form that travelled best at that moment. The table below follows one thread only: when the medium changes, the instrumentation changes with it.

\begin{longtable}{p{0.16\textwidth}p{0.34\textwidth}p{0.42\textwidth}}
\toprule
Period & Key shift & Effect on arrangement \\
\midrule
\endfirsthead
\toprule
Period & Key shift & Effect on arrangement \\
\midrule
\endhead
1900s--1920s & Tin Pan Alley song publishing, theatre and sheet music & Songs were designed to be sung and sold as sheet music; melody and lyric outranked recorded sound \\
1920s--1940s & Radio, records and sound film; Shanghai's shidaiqu and other Chinese pop industries & The microphone made quiet singing possible; small jazz bands and cabaret line-ups entered pop \\
1950s & Rock and roll, electric guitar, backbeat and the teenage market & Beats 2 and 4 became a body coordinate; the guitar rose from accompaniment to signature timbre \\
1960s & Motown, the British Invasion, the album and stadium shows & Tighter rhythm sections, thicker harmony; producers began to direct the sound like a film director \\
1970s & Singer-songwriters, soft rock, disco and synthesizers & One branch became more intimate (voice and acoustic guitar), another more dance-floor oriented (four-on-the-floor and string walls) \\
1980s & MTV, music video, synth-pop and drum machines & Visual memory entered songwriting; timbre itself became a hook, with drum machines and synths carrying the signature \\
1990s & R\&B, hip-hop, grunge, teen pop and global labels & Heavier low end and groove; sampling, rap and choral choruses ran in parallel \\
2000s & MP3, digital downloads, Auto-Tune and producer-centred records & Shorter intros, tighter loops; vocal tuning and stacking became standard craft \\
2010s & Streaming, EDM, global collaborations and short video & Songs had to survive repeated slicing; drops, loops and fifteen-second hooks entered form \\
2020s & Recommendation algorithms, short video, virtual production and genre mixing & Genre borders loosened; one song can borrow synthwave, afrobeat and hyperpop at once \\
\bottomrule
\end{longtable}

Chinese-language pop has its own lineage: Shanghai's 1930s shidaiqu joined film, records and cabaret; the Hong Kong and Taiwan industries of the 1980s--1990s built a full chain of idols, TV themes and cassette/CD releases; from the 2000s, R\&B, rock, folk and electronic vocabularies entered the mainstream. The industrial timing differs by region, but the arranging question is the same: how does a song get heard, remembered and sung along to in the medium of its time?

\section{Arrangement: one job per layer}
\label{sec:pop-craft}

Pop instrumentation can flex, but the division of labour must be clear. Writing ``who is responsible for what'' into the track names is more useful than auditioning ten presets first.

\subsection{Typical instrumentation and four layers}
\label{sec:pop-instrumentation}

\begin{longtable}{p{0.16\textwidth}p{0.25\textwidth}p{0.25\textwidth}p{0.20\textwidth}}
\toprule
Layer & Common instruments/tracks & Main job & Common mistake \\
\midrule
\endfirsthead
\toprule
Layer & Common instruments/tracks & Main job & Common mistake \\
\midrule
\endhead
Voice & Lead, double, harmony, ad-lib & Tell the story, carry the emotion, own the most important hook & Double sung in exactly the same pitch and rhythm, so it sounds like one person copied \\
Rhythm & Drums, bass, rhythm guitar/keys & Set tempo, set accents, wake the body & Filling every bar, so the listener cannot find a stable beat \\
Harmony & Pad, piano, chord guitar & Hold key and form, leave room for the voice & Pad written as a second melody, fighting the voice in the same register \\
Colour & Pluck, bell, brass, strings, FX & Mark section borders, create memory points, add ear candy & Colour running from top to tail, so the chorus has nothing fresh \\
\bottomrule
\end{longtable}

A practical starting point is eight tracks: lead vocal, double, drums, bass, pad, pluck, colour (brass/strings/bell), and the master bus. Run that skeleton first, then decide whether you need a second guitar, a second keyboard or more percussion. Track count is not proof of professionalism; every track having a job is.

\subsection{Groove: let the body understand first}
\label{sec:pop-groove}

The most common pop skeleton is 4/4 at 90--120 BPM. Kick on 1 and 3, snare on 2 and 4, closed hi-hat in eighths. This ``most stable skeleton'' is not laziness; it gives melody and lyric a predictable floor. Put the real change in:

\begin{itemize}
\item \textbf{the pre-chorus}: add sixteenth-note closed hats, lift the high end, or let the kick drop for half a beat;
\item \textbf{the chorus entrance}: land kick, bass root and one short colour accent on beat 1 together;
\item \textbf{the section border}: use a bar or two of fill or silence to refocus attention.
\end{itemize}

If every beat weighs the same, the chorus never ``arrives''. The rhythm section's craft is not ``something different in every bar''; it is knowing where to stay solid and where to change clearly.

\subsection{Harmony: put one nail into the familiar loop}
\label{sec:pop-harmony}

Pop harmony usually gives you a familiar loop plus one surprise. I--V--vi--IV and vi--IV--I--V are common because they are easy to sing, easy to remember and easy to place lyric on chord tones. Copying the loop exactly makes the song sound like an exercise. Without breaking singability, you can add one small nail:

\begin{itemize}
\item in the pre-chorus, replace the last chord with ii or IV/V to create suspension;
\item at the chorus entrance, borrow $\flat$VII or minor iv to turn the emotion darker or warmer;
\item in the verse, use an inverted bass so the same chords feel more intimate;
\item in the bridge, leave the key briefly, then return to the main loop in the outro.
\end{itemize}

Harmony is not better simply because it is more complex. In pop, the first job of a chord is to give the lyric somewhere to land and the section a direction; the second job is ``one surprise''. Too many surprises and the listener loses the safety to sing along.

\subsection{Melody: hook before sentence}
\label{sec:pop-melody}

The \term{topline}{topline} is the main melody at the top of the song, usually carried by the lead vocal. Write the smallest memorable unit first, then expand it into a sentence:

\begin{enumerate}
\item write a four-to-five-note motif that lands on chord tones;
\item give it a clear rhythmic outline (for example ``short--short--long'' or ``weak--strong'');
\item show only half of it in the verse, and the whole of it in the chorus;
\item use the pre-chorus to lift register or density so the first chorus note has the weight of arrival.
\end{enumerate}

Lyric and melody must make room for each other: important words sit on strong beats or long notes; connective words can be compressed, swallowed or placed on weak beats. The chorus hook often lands on beat 1 or 2. If you bury it on beat 3, the listener waits two extra beats; in a short-video era, that is expensive.

\subsection{Sound design: treat space as arrangement}
\label{sec:pop-sound}

Pop arrangement does not live only in the track list; it also lives in the space. \term{Ear candy}{ear candy} is the small sound that carries no main narrative but makes a section taste better: a bell, a reverse reverb, a filter sweep, a background harmony. It is salt in a dish: a little improves it, too much steals the scene.

Three executable principles:

\begin{itemize}
\item \textbf{Keep the verse dry}: keep the lead vocal forward, use little long reverb on pad and pluck, so the lyric stays clear.
\item \textbf{Give the chorus space}: the lead can take a short reverb/delay, the pad opens its high end, and the colour layer appears at the section entrance.
\item \textbf{One owner for the low end}: kick and bass take turns owning the bass; the pad gives its low end to them. If sidechain compression is available, use it; if not, draw the same ``give way'' curve on the pad's volume automation.
\end{itemize}

\section{Four songs, arrangement only}
\label{sec:pop-cases}

The four analyses below study arrangement only. They quote no lyrics and include no melodic notation. Work information follows album release information and official video pages; links are in the chapter's further reading.

\subsection{``Billie Jean'': the groove is the hook}
\label{sec:pop-case-billie-jean}

Michael Jackson's ``Billie Jean'' appears on \emph{Thriller} (1982) and was released as a single in 1983. What is most useful to learn is not ``how high the chorus is'', but ``how the whole song uses one groove as its hook''.

\begin{itemize}
\item \textbf{Rhythm}: the drum machine kick and snare are not complex; accents are clear and space is audible. The bass turns the low end into a memory point with a repeated eighth/sixteenth-note phrase.
\item \textbf{Voice}: the lead sits forward and fairly dry; there is breath between phrases. The chorus does not arrive through more tracks but through vocal weight and the entry of harmony.
\item \textbf{Colour}: synth strings hold long notes like a thin background; they do not fight the bass or cover the voice.
\item \textbf{Logic}: let the groove stand first, then use a small amount of colour to mark sections. There are not many tracks, but every one has a job.
\end{itemize}

Takeaway: if your song has no melody that gets hummed back, do not add tracks yet. Try making one small bass or drum motif into a loop and let it carry the hook.

\subsection{``Shape of You'': one loop can carry a song}
\label{sec:pop-case-shape}

Ed Sheeran's ``Shape of You'' was released in January 2017 as one of the lead singles from \emph{Divide}. It takes the loop to an extreme: a short marimba/pluck-like phrase repeats like a metronome and pulls the whole song along.

\begin{itemize}
\item \textbf{Loop}: the figure is not complicated, but its syncopation is clear and locks against the kick/clap with a half-beat offset.
\item \textbf{Sections}: the verse keeps space, letting the loop and bass speak; the chorus adds fuller low end, claps and vocal stacking while the loop remains the axis.
\item \textbf{Voice}: the melody prefers short motifs, repetition and variation over long sentences; strong lyric words land on the loop's accents.
\item \textbf{Logic}: it does not ``write new material for every section''; it changes density and space around the same material. That is why it survives repeated listening on streaming and short video.
\end{itemize}

Takeaway: build a four-bar loop first, then ask ``what do I remove in the verse and add back in the chorus?'' If the loop itself is not clear, more layers will not save it.

\subsection{``Blinding Lights'': retro is a filter, not a costume}
\label{sec:pop-case-blinding-lights}

The Weeknd's ``Blinding Lights'' was released in November 2019 and included on \emph{After Hours} in 2020. It borrows 1980s synth-pop timbres without turning into a retro cover: the low end, drum impact and vocal treatment are contemporary pop.

\begin{itemize}
\item \textbf{Timbre}: a synth arpeggio and wide string pad form the retro bed; the gated drum texture gives a clear pulse.
\item \textbf{Groove}: it feels like a half-time groove, but the pulses are dense; bass and kick hold the low end while the synth runs high.
\item \textbf{Voice}: the verse is restrained, the chorus opens; the vocal space changes with the section instead of using one reverb throughout.
\item \textbf{Logic}: the retro timbre is only the outer layer; the structure is still ``verse leaves space---chorus releases---tail collects''.
\end{itemize}

Takeaway: to use retro timbres, do not just change presets. Decide first which layer carries the pulse, which carries width and which carries emotion; then let the timbres serve those three jobs.

\subsection{``bad guy'': space and low end as the hook}
\label{sec:pop-case-bad-guy}

Billie Eilish's ``bad guy'' was released in March 2019 on \emph{When We All Fall Asleep, Where Do We Go?}. Its hook is not a high melodic phrase but an overall texture: low, empty and strange.

\begin{itemize}
\item \textbf{Low end}: the bass stays low, and kick and claps do not fill every gap; the weight of the low end is itself the memory point.
\item \textbf{Space}: the voice is light and close, with obvious emptiness between phrases; some effects seem to come from a corner of the room.
\item \textbf{Contrast}: sections switch density and register suddenly instead of stacking layer upon layer; sometimes ``one beat less'' is stronger than ``one layer more''.
\item \textbf{Logic}: timbre, space and dynamics make the hook, so the listener remembers the song's ``smell'', not only one melodic phrase.
\end{itemize}

Takeaway: if your voice or melody is not a strong hook, design a distinctive low-end/timbre/space combination first; but do not be strange for the sake of strange. Every timbre must serve the emotion.

\section{The yeban skeleton: eight bars of pop}
\label{sec:pop-yeban}

The skeleton below can be rebuilt directly in yeban. It is not a complete song; it is a run-through of the ``verse--chorus'' arranging logic.

\begin{longtable}{p{0.10\textwidth}p{0.22\textwidth}p{0.58\textwidth}}
\toprule
Bar & Section & Arranging action \\
\midrule
\endfirsthead
\toprule
Bar & Section & Arranging action \\
\midrule
\endhead
1--4 & Verse & Drums: kick 1/3, snare 2/4, closed hats in eighths; bass: root on 1/3; pad: long notes; pluck: only the last two beats of a phrase; lead: middle-low register \\
5--8 & Chorus & Drums: add open hat/clap, short fill on bar 5 beat 1; bass: root plus octave; pad: open the high end; pluck: switch to eighth-note arpeggio; lead: add a low-third double; short brass/bell accent on bar 5 beat 1 \\
\bottomrule
\end{longtable}

Order of operations:

\begin{sloppypar}
\begin{enumerate}
\item Create a 90 BPM, 4/4 project; build and name the tracks above. Use audio tracks for lead/double, MIDI for drums/bass/pad/pluck/colour, and keep the master as a bus.
\item Write drums and bass first. Solo them, confirm the low end has one owner, then open the mixer with F8 and check that the pad's low end is not pushing into the bass.
\item Write the pad's long notes and the pluck's phrase endings. In the verse, let the pluck show only two beats, not wall-to-wall.
\item Write the lead: target chord tones C/G/Am/F; leave half a beat of breath at the end of bar 4 so the chorus downbeat on bar 5 has somewhere to land.
\item On bar 5, add the double, open hat and colour accent. Add one layer only; do not pour every piece of ear candy in at once.
\item Use F7 for the piano roll to fix notes and F3 for the stacked view to compare instruments and effects. Solo after every added layer and confirm it is helping.
\item Use Ctrl+Shift+A for automation: draw a high-end opening curve for the pad in the chorus. If sidechain is unavailable, draw the give-way curve on the pad's volume instead.
\item Save the project and export an eight-bar audition. If removing a track makes the music worse, keep it; if you cannot answer, delete it or push it into the background.
\end{enumerate}
\end{sloppypar}

\section{Exercises: three paths, different trade-offs}
\label{sec:pop-practice}

\begin{exercise}
\textbf{An eight-bar loop, three starting paths.} Using the same C--G--Am--F, make three versions: A copy the groove of a reference track; B hum a phrase first and then notate it; C chop your own recording into a loop. Compare ``identity, speed of completion, editability''. A is fast but tends to sound like the original; B is slower but more yours; C has texture but carries sample clearance and cleanup costs.
\end{exercise}

\begin{exercise}
\textbf{Three ways to lift a pre-chorus.} Attach three different pre-choruses to the same verse: A lift by harmony (IV/V or ii); B subtract arrangement (drop the drums, then return); C rhythmic fill (sixteenth hats plus a snare roll). Play each and note which one spends the chorus's energy too early; try all three before choosing what your song needs.
\end{exercise}

\begin{exercise}
\textbf{A five-note hook.} Limit yourself to five notes for a chorus motif and forbid long sentences. Write three versions: one all on strong beats, one with the accent displaced to a weak beat, one with the last note left empty. Ask two people to listen blind and say which version they can still hum ten minutes later. Test memory through melody and rhythm only, without lyric prompts.
\end{exercise}

\section{Further reading}
\label{sec:pop-reading}

\begin{sloppypar}
\begin{itemize}
\item \textbf{AllMusic, ``Pop/Rock''}: \url{https://www.allmusic.com/genre/pop-rock-ma0000002613}. Use the genre entry to fill in the pop/rock lineage and representative artists; good for building a coordinate system first.
\item \textbf{Michael Jackson, ``Billie Jean'' (Official Video)}: \url{https://www.youtube.com/watch?v=Zi_XLOBDo_Y}. Watch how groove, space and voice create a hook without stacking tracks.
\item \textbf{Ed Sheeran, ``Shape of You'' (Official Music Video)}: \url{https://www.youtube.com/watch?v=JGwWNGJdvx8}. Hear how one short loop carries the structure through density changes between verse and chorus.
\item \textbf{The Weeknd, ``Blinding Lights'' (Official Video)}: \url{https://www.youtube.com/watch?v=4NRXx6U8ABQ}. Compare the retro synth palette with its contemporary low end and drum impact.
\item \textbf{Billie Eilish, ``bad guy'' (Official Video)}: \url{https://www.youtube.com/watch?v=DyDfgMOUjCI}. Observe how low end, space and dynamic contrast replace a traditional high-note hook.
\end{itemize}
\end{sloppypar}

The links above were checked one by one against official pages and search results for title and author on 27 November 2026. If a platform's regional restriction prevents opening one, search for the entry name on the artist's official account in that platform.

\section{Chapter takeaways}
\label{sec:pop-takeaways}

\begin{itemize}
\item Pop is not ``simple''; it is the clear handling of hook, groove, contrast and payoff.
\item Genre evolution follows the medium: sheet music, radio, records, MTV, streaming and short video each changed the priority of arrangement.
\item Write function before instrument; one job per layer; track count is not proof of professionalism.
\item Harmony uses a familiar loop for safety, then one suspension or borrowed chord for memory.
\item The topline starts with a minimal motif and expands into a sentence; put the chorus hook in the first two beats where possible.
\item The yeban skeleton runs ``drums/bass---harmony---lead---colour---automation---solo check'' in that order.
\item Exercises should compare the trade-offs of different paths, not just chase a reference track; editable, removable and explainable is usable arrangement.
\end{itemize}

\begin{helpmapping}
\begin{longtable}{p{0.45\textwidth}p{0.45\textwidth}}
\toprule
Section & Help anchor \\
\midrule
What pop is & \texttt{arrangement-genre-pop-what} \\
From gramophone to playlist & \texttt{arrangement-genre-pop-history} \\
Arrangement: one job per layer & \texttt{arrangement-genre-pop-craft} \\
Four songs, arrangement only & \texttt{arrangement-genre-pop-cases} \\
The yeban skeleton & \texttt{arrangement-genre-pop-yeban} \\
Exercises & \texttt{arrangement-genre-pop-practice} \\
Further reading & \texttt{arrangement-genre-pop-reading} \\
\bottomrule
\end{longtable}
\end{helpmapping}
